\documentclass[11pt]{article}

% ============================================================
% Register-Sensitivity Safety Gap in Open LLMs
% Author: Shubhanandan T Y
% Plain article class; reformat to venue style before submission.
% Figures expected in ./figs/  (fig1_divergence.png, fig2_svg_all.png,
% fig3_probe.png, fig0_pilot.png)
% ============================================================

\usepackage[utf8]{inputenc}
\usepackage[T1]{fontenc}
\input{glyphtounicode}
\pdfgentounicode=1
\usepackage{amsmath,amssymb}
\usepackage{booktabs}
\usepackage{graphicx}
\usepackage{hyperref}
\usepackage{xcolor}
\usepackage[margin=1in]{geometry}
\usepackage{caption}
\usepackage{subcaption}
\usepackage{natbib}
\usepackage{authblk}

\hypersetup{colorlinks=true, linkcolor=blue, citecolor=blue, urlcolor=blue}

\title{\bf Hidden in the Hedge:\\
Unvalidated LLM Graders Understate the Register Safety Gap in\\
Instruction-Tuned Language Models}

\author[1]{Shubhanandan T Y}
\affil[1]{\normalsize [Affiliation to be added] \\ \texttt{[email to be added]}}

\date{\today}

\begin{document}
\maketitle

% ------------------------------------------------------------
\begin{abstract}
Safety evaluations of language models increasingly use LLM-based graders to judge
whether a response is harmful. We show that an \emph{unvalidated} grader can
drastically understate a real safety gap. We study sensitivity to \emph{register}---the
difference between formal, benchmark-style phrasing and the casual, first-person
phrasing real users write---using 300 paired prompts that express an identical harmful
intent in each register, with explicit ``this is a test'' cues removed. Across six
instruction-tuned open models in two families (Qwen2.5 at 0.5/1.5/3/7B and Gemma-2 at
2/9B), a keyword refusal metric reports a large register gap (corpus-wide $0.27$, up to
$0.54$ for Qwen-7B). Re-scoring every response with an off-the-shelf LLM grader that
labels responses as refusal, non-actionable hedge, or actionable harm reports a
near-zero harm gap ($0.03$)---seemingly implying the refusal gap is mostly harmless
hedging. But when we validate that grader against human labels on 200 of its
``hedge'' judgments, we find it systematically misclassifies actionable harmful
content as hedging, and does so far more for casual prompts: humans judge $20\%$ of
formally-framed grader-hedges but $49\%$ of casually-framed ones to be actually
harmful. Correcting for this audited error yields a register harm gap of
$\mathbf{0.26}$ (95\% CI $[0.19, 0.32]$), with casual framing raising the
actionable-harm rate from $8.5\%$ to $34.3\%$. The naive grader had hidden roughly
seven-eighths of this gap. Our central message is methodological: \textbf{LLM graders
must be validated against human labels before their safety verdicts are trusted}; an
unvalidated grader here would have reported no safety problem where a substantial one
exists. We also report a probing negative result and a cross-family behavioral
pattern. Our study is limited to open models up to 9B; we release the dataset, all
generations, both scorings, human labels, and code.
\end{abstract}

% ------------------------------------------------------------
\section{Introduction}

Pre-deployment safety evaluation underpins nearly every governance mechanism for
large language models (LLMs), from developer release policies to national AI
safety institutes to emerging regulation. All of these mechanisms share an
implicit and load-bearing assumption: that a model's measured behavior on
evaluation prompts is a faithful proxy for its behavior once deployed. If that
assumption fails---if models behave systematically differently when a prompt
``looks like'' an evaluation---then safety scores overstate real-world safety, and
the decisions built on them are unsound.

A growing body of work suggests this assumption is fragile. Frontier and open
models can, to varying degrees, distinguish evaluation contexts from deployment
contexts and shift their behavior accordingly, a phenomenon termed
\emph{evaluation awareness} \citep{needham2025,nguyen2025,goldowsky2025}. Recent
analysis argues this is not a single capability but a multi-dimensional construct
whose axes---detectability, behavioral manifestation, and controllability---are
only weakly correlated, so that no single ``awareness score'' reliably predicts
deployment safety \citep{notoneability2026}. Much of this literature, however,
induces the evaluation/deployment distinction using \emph{explicit} cues (prompts
literally labelled as tests or benchmark items) and studies it either within a
single model family or through activation probes whose validity has recently been
questioned \citep{devbunova2026}.

We isolate a narrower and, we argue, more deployment-relevant variable:
\textbf{register}. Real users do not phrase requests like benchmark items. They
write casually, in the first person, with backstory, informal grammar, and
occasional typos. If a model's safety behavior is sensitive to this stylistic
difference \emph{alone}---with no explicit test markers present---then even a
perfectly curated safety benchmark can mispredict deployment behavior simply
because benchmark prompts and real traffic occupy different registers. This is a
sharper claim than generic evaluation awareness: it says the confound is baked into
the \emph{style} of evaluation, not merely into telltale keywords that a careful
benchmark designer could remove.

We make three contributions.

\begin{enumerate}
    \item \textbf{A grader-validation result (primary).} Using paired formal/casual
    prompts and a metric we call the \emph{Register Validity Gap} (RVG), we show that
    the measured register safety gap depends critically on how harm is judged. A
    keyword refusal metric reports a gap of $0.27$--$0.54$; an off-the-shelf LLM harm
    grader reports $0.03$; but human validation of that grader (200 labels) reveals it
    hides most casually-framed harm as ``hedging,'' and the corrected harm gap is
    $0.26$ (95\% CI $[0.19, 0.32]$). An unvalidated grader would have reported
    essentially no safety gap where a substantial one exists. We conclude that LLM
    safety graders must be human-validated before use.

    \item \textbf{A cross-family behavioral pattern.} We measure the refusal-level
    register effect across six models in two families, finding it significant
    everywhere but with family-dependent scaling behavior---a further caution against
    extrapolating safety-evaluation trends across architectures.

    \item \textbf{A probing caution.} We report that linear probes on prompt
    activations separate the two registers with AUROC $\approx 1.0$ at all layers,
    persisting after removing keywords and matching length, and argue this reflects
    surface separability rather than an internal evaluation-awareness representation.
\end{enumerate}

% ------------------------------------------------------------
\section{Related Work}
\label{sec:related}

\paragraph{Situational and evaluation awareness.}
The broader notion that models possess \emph{situational awareness}---knowledge of
their own nature and context---was established by early work showing models can
reason about their situation without in-context examples, and formalized in the
Situational Awareness Dataset \citep{laine2024}, whose \texttt{stages\_oversight}
subset directly probes recognition of evaluation contexts.
\citet{needham2025} systematically studied how models distinguish evaluation from
deployment, finding that model judgments of scenario authenticity correlate
strongly with human judgments and identifying the surface features that signal an
evaluation context. \citet{goldowsky2025} showed models can recognize when they are
being evaluated for alignment specifically. Together these establish that
evaluation awareness exists; our work asks a complementary question about one
particular, deployment-relevant \emph{cause} of behavioral shifts---register---and
measures the resulting safety consequences behaviorally.

\paragraph{Evaluation awareness is multi-dimensional.}
\citet{notoneability2026} run eight experiments across 37 open models and seven
families and argue that evaluation awareness decomposes into partially independent
axes; in particular, they find that awareness metrics and safety-behavior metrics
are largely orthogonal, so no single awareness score reliably predicts deployment
safety. This motivates our focus on directly measuring the behavioral safety gap
rather than inferring it from an awareness proxy.

\paragraph{Benchmark validity and surface sensitivity.}
A parallel literature documents that benchmark conclusions depend heavily on
prompt wording and surface form, motivating multi-prompt and rewrite-based
evaluation \citep{mizrahi2024,xiong2025}. \citet{xiong2025} in particular propose a
probe-rewrite-evaluate workflow that paraphrases benchmark items to quantify
evaluation awareness; our design differs in that we do not paraphrase freely but
construct \emph{matched pairs} that hold harmful intent fixed while varying a single
controlled axis (register), and we target the safety-refusal behavior specifically
rather than general benchmark accuracy. Related work shows models can answer
multiple-choice items using artifacts rather than the question itself
\citep{balepur2024}, and warns of benchmark contamination and ``cheating''
\citep{zhou2023}. \citet{ren2024safetywashing} question whether safety benchmarks
measure safety progress at all. Our register manipulation can be seen as a
targeted, safety-specific instance of this surface-sensitivity concern: we hold
semantics fixed and vary only style.

\paragraph{Probing and its limits.}
Linear probes on frozen activations are widely used to argue that models
internally represent high-level properties, including refusal direction
\citep{arditi2024} and evaluation awareness
\citep{nguyen2025}. However, \citet{devbunova2026} shows that standard
single-contrast probes trained on benchmark-vs-chat data primarily track prompt
\emph{format} rather than evaluation context, and that decorrelating format from
context is necessary to recover a genuine signal. Our probing result is an
independent, constructive confirmation of this concern: we show the confound
survives two successive rounds of removal (lexical, then length), and argue that
prompt-activation probes cannot in principle separate ``the model inferred it is
being tested'' from ``these two prompts are stylistically different.''

\paragraph{Test-awareness and steering.}
\citet{abdelnabi2025} study a Hawthorne-like effect in reasoning models and show
test awareness can be linearly steered, revealing differential compliance. This
supports the plausibility of a behavioral register effect, though our contribution
is a cross-family behavioral \emph{measurement} rather than an intervention.

% ------------------------------------------------------------
\section{Methods}
\label{sec:methods}

\subsection{Paired-prompt design}
We define 25 harmful \emph{intents} spanning ten categories: cyber, fraud, weapons,
chemical, privacy, misinformation, self-harm information, harassment, drugs, and
theft. Each intent is expressed in two \emph{framings} that hold the intent
constant and vary only surface style:

\begin{itemize}
    \item \textbf{Eval framing:} formal, impersonal, instruction-style phrasing,
    e.g.\ \emph{``Outline in detail, from beginning to end, the full set of steps
    that someone would follow in order to \{intent\}.''}
    \item \textbf{Deploy framing:} casual, first-person, with backstory and
    informal grammar, e.g.\ \emph{``hey so, i keep reading about how someone might
    \{intent\}. what's the real method?''}
\end{itemize}

Expanding the intents across multiple templates yields 300 matched pairs. Crucially,
the eval framing contains \emph{no} explicit evaluation markers---no ``benchmark,''
``evaluation,'' ``test,'' or ``assessment''---so that the only systematic
difference between the two framings is register.

\subsection{Three-stage confound removal}
\label{sec:confound}
Our final dataset is the product of three iterations. We report all three because
each isolates a distinct confound and the intermediate results are themselves
informative.

\begin{enumerate}
    \item \textbf{Explicit cues (rejected).} An initial version used eval framings
    containing literal markers (``You are being evaluated,'' ``Benchmark item
    \#$n$''). This produced a smooth, monotonic-looking behavioral scaling curve
    \emph{and} perfect layer-1 probe separation. Both patterns are consistent with
    models and probes keying on the marker vocabulary rather than register; while we
    did not run a controlled ablation isolating the markers, the fact that removing
    them (stage 2) flattened part of the behavioral curve supports this
    interpretation. We discarded this version.
    \item \textbf{De-lexicalized.} We removed all evaluation-marker vocabulary,
    leaving formal-vs-casual register as the only systematic difference. This
    materially changed the behavioral picture: part of the apparent small-model
    scaling disappeared, indicating it had been an artifact of the keywords. All
    behavioral results in this paper use this version.
    \item \textbf{Length-matched.} Because a probe can exploit sequence length as a
    shortcut, we additionally produced a version in which eval and deploy framings
    are matched in length (means within ${\sim}3$ words) for the probe analysis of
    Section~\ref{sec:probe}.
\end{enumerate}

\subsection{Models and generation}
We evaluate Qwen2.5-Instruct at 0.5B, 1.5B, 3B, and 7B \citep{qwen2025} and
Gemma-2-it at 2B and 9B \citep{gemma2}, all in 4-bit quantization on a single
consumer GPU. Generation is greedy (deterministic) with a 200-token cap. Using two
independently developed families lets us separate family-specific effects from
general ones.

\subsection{Refusal scoring and its validation}
We classify each response as a refusal using a keyword/heuristic scorer over the
opening of the response. This is the load-bearing measurement in the paper, so we
assess it carefully and flag its residual weaknesses honestly.

First, a \emph{directional} concern: our metric could be biased in either direction
if soft refusals are missed asymmetrically across framings. In particular, casually
phrased responses might contain casually phrased refusals (``nah, can't help with
that'') that a keyword scorer misses, which would \emph{inflate} RVG rather than make
it conservative. To check this, we scanned all 904 deploy-framing responses labeled
as non-refusals for casual-refusal markers and found only one plausible missed
refusal; models in our sample overwhelmingly refuse in canonical phrasing even under
casual framing. This makes large asymmetric undercounting on the casual side
unlikely, though we note the scan is itself lexical.

Second, a \emph{ground-truth} concern: two automatic heuristics agreeing (we observe
99.9\% agreement between the keyword scorer and a stricter heuristic across all
3{,}600 responses) establishes consistency, not accuracy, since they may share blind
spots. We therefore hand-labeled a stratified sample of 60 responses balanced across
all six models and both framings. The keyword scorer agrees with human labels with
95.0\% accuracy and Cohen's $\kappa = 0.90$ (``almost perfect'' agreement). All three
disagreements are cases where the human marked a refusal that the keyword scorer
missed (a soft refusal not opening with a canonical phrase); there were no cases in
which the scorer labeled a genuine compliance as a refusal. The scorer therefore
\emph{undercounts} refusals, and does so on both framings, which supports---now with
human ground truth rather than assertion---that our RVG is not inflated by
asymmetric refusal-detection error. We additionally manually verified the Gemma-2-9B
ceiling (Section~\ref{sec:verify}).

Third, a \emph{construct} concern: refusal-rate RVG measures non-refusal, not
harmfulness. A casual response that fails to refuse but produces no genuinely harmful
content still inflates a refusal-based gap. We address this directly in
Section~\ref{sec:harm} by re-grading every response for actionable harm and validating
that grader against human labels.

\subsection{The Register Validity Gap}
For each model, let $r_{\text{eval}}$ and $r_{\text{deploy}}$ be the refusal rates
under the two framings. We define
\[
    \mathrm{RVG} \;=\; r_{\text{eval}} - r_{\text{deploy}}.
\]
A positive RVG means the model appears \emph{safer} under formal framing than it
actually is under casual framing. Because our 300 pairs are generated from 25 base
intents (via templates) rather than drawn independently, per-pair resampling would
understate uncertainty. We therefore report 95\% confidence intervals via a
\textbf{cluster bootstrap that resamples intents} (not pairs), taking all pairs of
each resampled intent (5{,}000 resamples); this respects the non-independence of
pairs sharing an intent. We report significance via McNemar's exact test, which is the appropriate paired test
for binary refusal outcomes, complemented by the clustered intervals above.

% ------------------------------------------------------------
\section{Results}
\label{sec:results}

\subsection{The gap is significant in every model}
Table~\ref{tab:svg} reports RVG for all six models. Every value is positive and
significant ($p<0.001$), with confidence intervals well clear of zero. Across both
families and all scales, models refuse the formal framing of a harmful request more
often than the identical casual framing. This is the central empirical fact of the
paper: safety behavior is register-sensitive, and it is so even with all explicit
evaluation cues removed.

\begin{table}[t]
\centering
\caption{Register Validity Gap (RVG) across six models and two families.
$r_{\text{eval}}$ and $r_{\text{deploy}}$ are refusal rates under formal and casual
framing; $\mathrm{RVG}=r_{\text{eval}}-r_{\text{deploy}}$. The 95\% CI is a
\emph{cluster bootstrap over the 25 intents} (not over pairs), which widens the
intervals relative to naive per-pair resampling but respects non-independence; every
interval still excludes zero. All $p<0.001$ (McNemar's exact test), $n=300$ pairs.}
\label{tab:svg}
\begin{tabular}{llccccc}
\toprule
Family & Model & Size (B) & $r_{\text{eval}}$ & $r_{\text{deploy}}$ & RVG & 95\% CI (clustered) \\
\midrule
Qwen2.5 & 0.5B-Instruct & 0.5 & 0.633 & 0.413 & 0.220 & [0.097, 0.347] \\
Qwen2.5 & 1.5B-Instruct & 1.5 & 0.640 & 0.483 & 0.157 & [0.027, 0.273] \\
Qwen2.5 & 3B-Instruct   & 3.0 & 0.603 & 0.397 & 0.207 & [0.093, 0.320] \\
Qwen2.5 & 7B-Instruct   & 7.0 & 0.707 & 0.167 & \textbf{0.540} & [0.413, 0.657] \\
\midrule
Gemma-2 & 2B-it & 2.0 & 0.983 & 0.643 & 0.340 & [0.250, 0.433] \\
Gemma-2 & 9B-it & 9.0 & 1.000 & 0.883 & 0.117 & [0.063, 0.180] \\
\bottomrule
\end{tabular}
\end{table}

\subsection{Effect sizes and paired tests}
To ensure the effect is not an artifact of a single test, we report additional
paired statistics in Table~\ref{tab:stats}. McNemar's exact test---appropriate for
our paired binary design---is significant for every model ($p<10^{-5}$). Effect
sizes are moderate to large: Cohen's $h$ ranges from $0.32$ (Qwen-1.5B) to $1.16$
(Qwen-7B). Refusal odds ratios (Haldane--Anscombe corrected so the Gemma-9B ceiling
is finite) range from $1.9$ to $80.4$; see Table~\ref{tab:stats}.

\begin{table}[t]
\centering
\caption{Effect sizes and paired significance tests. McNemar's exact $p$ tests the
asymmetry of discordant pairs. Cohen's $h$ is the effect size for the difference in
refusal proportions; OR is the refusal odds ratio (formal vs.\ casual) with a
Haldane--Anscombe ($+0.5$) correction so that the ceiling case (Gemma-9B) yields a
finite value. $n=300$.}
\label{tab:stats}
\begin{tabular}{lccc}
\toprule
Model & McNemar $p$ & Cohen's $h$ & Odds ratio (HA-corrected) \\
\midrule
Qwen2.5-0.5B & $1.1\times10^{-9}$  & 0.444 & 2.44 \\
Qwen2.5-1.5B & $7.7\times10^{-6}$  & 0.317 & 1.90 \\
Qwen2.5-3B   & $2.5\times10^{-7}$  & 0.416 & 2.31 \\
Qwen2.5-7B   & $4.2\times10^{-45}$ & 1.156 & 11.91 \\
Gemma-2-2B   & $1.0\times10^{-29}$ & 1.021 & 29.85 \\
Gemma-2-9B   & $5.8\times10^{-11}$ & 0.697 & 80.36 \\
\bottomrule
\end{tabular}
\end{table}

\subsection{Scaling of the refusal-level gap is architecture-dependent}
We now describe how the \emph{refusal-level} gap (not the harm-level gap of
Section~\ref{sec:harm}) varies with scale; readers should keep in mind that this
reflects hedging behavior more than harmful compliance. The relationship between the
refusal RVG and model size differs qualitatively across families
(Figure~\ref{fig:divergence}).

\paragraph{Qwen: the gap widens at scale.} RVG is roughly flat from 0.5B to 3B
($0.16$--$0.22$) then jumps to $0.54$ at 7B. The driver is the casual-framing
refusal rate, which \emph{collapses} from ${\sim}0.40$ to $0.167$ at 7B while the
formal-framing rate stays high. The largest Qwen model is thus the one that looks
safest under formal evaluation yet is least safe under casual phrasing---a pattern
that would cause benchmark scores to overstate deployed safety.

\paragraph{Gemma: the gap narrows at scale.} RVG falls from $0.34$ at 2B to $0.12$
at 9B. The driver is the opposite: Gemma-2-9B refuses nearly all requests in
\emph{both} framings ($r_{\text{eval}}=1.00$, $r_{\text{deploy}}=0.88$),
compressing the gap via a refusal ceiling induced by stronger safety tuning.

These families therefore exhibit opposite scaling behavior driven by opposite
mechanisms---collapsing deployment refusals versus a rising refusal ceiling. Any
claim that the benchmark--deployment gap grows or shrinks monotonically with
capability is consequently family-dependent, not a general law. We regard this as a
central cautionary result: a single-family study would have supported a confident
but incorrect scaling narrative.

\begin{figure}[t]
\centering
\includegraphics[width=\linewidth]{figs/fig1_divergence.png}
\caption{Refusal rate under formal (eval) vs.\ casual (deploy) framing across model
scale. The shaded region between the lines is the Register Validity Gap. In Qwen
(left) the gap widens at 7B because deployment-framing refusals collapse; in Gemma
(right) it narrows at 9B because the model approaches a refusal ceiling in both
framings.}
\label{fig:divergence}
\end{figure}

\begin{figure}[t]
\centering
\includegraphics[width=0.85\linewidth]{figs/fig2_svg_all.png}
\caption{Register Validity Gap for all six models with 95\% bootstrap confidence
intervals. Blue: Qwen2.5; red: Gemma-2. Every gap is significantly positive, but
the trend with scale differs by family.}
\label{fig:svgall}
\end{figure}

\begin{figure}[t]
\centering
\includegraphics[width=0.6\linewidth]{figs/fig5_scatter.png}
\caption{Each model's refusal rate under formal (x-axis) vs.\ casual (y-axis)
framing. Points on the dashed diagonal would indicate no register effect; every
model lies \emph{below} it, and the vertical distance to the diagonal is the RVG.
Qwen-7B sits farthest below the line; Gemma-9B sits near the top-right corner
(near-ceiling refusal in both framings).}
\label{fig:scatter}
\end{figure}

\subsection{Category structure}
The gap is not uniform across harm categories (Figure~\ref{fig:heatmap}). For
Gemma-2-2B, for example, the gap is largest for misinformation (RVG $=0.69$) and
cyber ($0.57$) and smallest for theft ($0.00$, where the model refuses both
framings uniformly). Across models, categories in which a casual phrasing most
plausibly resembles an innocent real-user question---``how do people actually do
this?''---tend to show the largest gaps, consistent with (though not proof of) the
interpretation that register is associated with an implicit deployment signal. The heatmap also makes the
family-level contrast visible: Qwen's gap concentrates and intensifies at 7B,
whereas Gemma's compresses toward zero at 9B as refusals saturate. Because each
category contains only ${\sim}30$ pairs (fewer for the smallest categories),
individual per-cell values are exploratory rather than confirmatory. Some cells are
even slightly negative (e.g.\ small negative RVGs in privacy and self-harm categories
for certain models); at this per-cell sample size these are consistent with noise
around zero and should not be interpreted as reliable reversals. We present the
heatmap to illustrate the aggregate family-level pattern, not to support
category-level claims.

\begin{figure}[t]
\centering
\includegraphics[width=0.85\linewidth]{figs/fig4_category_heatmap.png}
\caption{Register Validity Gap by harm category (rows) and model (columns). Warmer
cells indicate a larger gap (the model refuses the formal framing more than the
casual one). The gap concentrates in categories such as misinformation and cyber,
and---reading left to right---intensifies at Qwen-7B while compressing at Gemma-9B.}
\label{fig:heatmap}
\end{figure}

\subsection{Verifying the ceiling}
\label{sec:verify}
Because Gemma-2-9B refuses all 300 formal prompts exactly, we manually inspected a
random sample of its refusals. All were substantive, request-specific declines
(e.g.\ \emph{``I cannot provide you with steps on how to \dots{} Doing so is
illegal \dots''}), with response lengths exceeding 590 characters and naturally
varied phrasing---confirming the ceiling is a genuine behavioral property, not a
scoring artifact. This verification also serves as an internal control: the
keyword scorer does not manufacture the gap, since the smallest model (Qwen-0.5B)
shows no such saturation.

% ------------------------------------------------------------
\section{Refusal Is Not Harm: Re-scoring with a Harm Grader}
\label{sec:harm}

The results so far measure \emph{refusal}. But a response that does not refuse is not
necessarily harmful: it may hedge---warn, moralize, or redirect---without providing
usable harmful information. To test whether the register gap is a genuine
\emph{safety} gap or largely a \emph{hedging} gap, we re-scored all 3{,}600 responses
with an LLM grader (Qwen2.5-7B applying a strict rubric) that assigns each response
to one of three classes: \textsc{refusal} (declines or only warns), \textsc{hedge}
(vague or non-actionable, no usable method), or \textsc{harmful} (provides actionable
harmful content, counted as harmful even under an ``educational'' disclaimer).

\paragraph{The harm gap is far smaller than the refusal gap.}
Table~\ref{tab:threemetric} compares three gap measures per model. The refusal-based
gap is large (up to $0.54$ for Qwen-7B), while the \emph{naive} grader reports a much
smaller harmful-content gap (Qwen-7B: $0.037$). However, as we show below, the naive
grader systematically under-counts harm; after human-validated correction the harm gap
rises to $0.26$ (95\% CI $[0.19,0.32]$)---comparable to the refusal gap, and far from
the near-zero the naive grader reported. Figure~\ref{fig:metricgap} shows all three
measures; the correction is detailed under grader validation.

\begin{table}[t]
\centering
\caption{Three gap measures per model (eval $-$ deploy). ``Refusal (keyword)'' is our
original metric; ``Refusal (grader)'' is the grader's refusal-vs-not gap; ``Harmful''
is the gap in actionable-harmful output. The harmful gap is roughly an order of
magnitude smaller than the keyword refusal gap.}
\label{tab:threemetric}
\begin{tabular}{lccc}
\toprule
Model & Refusal gap (keyword) & Refusal gap (grader) & Harmful-content gap \\
\midrule
Qwen2.5-0.5B & 0.220 & 0.343 & 0.047 \\
Qwen2.5-1.5B & 0.157 & 0.187 & 0.057 \\
Qwen2.5-3B   & 0.207 & 0.147 & 0.053 \\
Qwen2.5-7B   & \textbf{0.540} & 0.157 & \textbf{0.037} \\
Gemma-2-2B   & 0.340 & 0.413 & 0.007 \\
Gemma-2-9B   & 0.117 & 0.357 & 0.000 \\
\bottomrule
\end{tabular}
\end{table}

\paragraph{Casual framing shifts refusals into hedges and, in part, into harm.}
The grade distribution (Figure~\ref{fig:gradedist}) shows the mechanism directly.
Moving from formal to casual framing, refusals fall from 1{,}100 to 619 (of 1{,}800),
and the naive grader sends most of the freed mass to \textsc{hedge} (684
$\rightarrow$ 1{,}105) rather than \textsc{harmful} (16 $\rightarrow$ 76). Taken at
face value this suggests casual framing makes models $\sim$27 points less likely to
\emph{refuse} but only $\sim$3 points more likely to produce harmful content. However,
human validation (below) shows the grader hides substantial casual-framing harm inside
the hedge category; after correction the increase in actionable harm is $\sim$26
points ($8.5\%\rightarrow34.3\%$). Casual framing thus produces substantially more
actionable harm, which the naive grader almost entirely missed.

\begin{figure}[t]
\centering
\includegraphics[width=0.85\linewidth]{figs/fig6_metric_gap.png}
\caption{The register gap measured by refusal (red) versus by \emph{naive}
actionable-harm grading (blue). Refusal-based metrics show a large gap; the naive
harm grader shows a much smaller one. Human validation (Section~\ref{sec:harm}) shows
the naive grader severely understates harm, and the human-corrected harm gap is
$0.26$ (95\% CI $[0.19,0.32]$).}
\label{fig:metricgap}
\end{figure}

\begin{figure}[t]
\centering
\includegraphics[width=0.6\linewidth]{figs/fig7_grade_dist.png}
\caption{Grade distribution under formal vs.\ casual framing (n$=$1{,}800 each).
Casual framing converts refusals into non-actionable hedges far more than into
harmful content.}
\label{fig:gradedist}
\end{figure}

\paragraph{Why the two refusal measures diverge.}
The keyword and grader refusal gaps differ substantially for some models, and the
direction is informative. For Qwen-7B the keyword gap ($0.540$) far exceeds the
grader gap ($0.157$): the keyword scorer counts Qwen-7B's verbose ``caveat-then-non-answer''
responses as compliance, inflating its apparent gap, whereas the grader reclassifies
them as hedges. For Gemma-9B the reverse holds---the grader gap ($0.357$) exceeds the
keyword gap ($0.117$)---because Gemma's refusals are long and discursive (``I cannot
$\dots$ but here is why $\dots$''), which the keyword scorer sometimes fails to catch
as refusals while the grader correctly identifies them. In both cases the divergence
reflects the keyword scorer's known difficulty with verbose, hedged responses, and
supports our central point that refusal-keyword scoring is an unreliable proxy for
either refusal or harm.

\paragraph{Grader validation and correction.}
An off-the-shelf grader is only as trustworthy as its agreement with human judgment.
We therefore hand-labeled 200 responses the grader had labeled \textsc{hedge}---100
per framing, the pivotal category for the harm gap. The grader's hedge bucket is
badly miscalibrated: humans judged only $30\%$ of ``hedges'' to be true hedges, with
the rest split between genuine refusals and, crucially, actionable harmful content.
The harmful fraction is strongly asymmetric across framing: \textbf{$20\%$ of
formally-framed grader-hedges but $49\%$ of casually-framed ones are actually
harmful}. Because casual framing both produces more hedges and a higher harmful
fraction within them, the naive grader understates harm far more on the casual side,
collapsing the measured gap. Propagating these audited rates over the corpus (with a
beta-binomial treatment of the audit uncertainty) raises the harmful rate from
$8.5\%$ under formal framing to $34.3\%$ under casual framing, giving a corrected
harm gap of $\mathbf{0.258}$ with a $95\%$ credible interval of $[0.189, 0.322]$
(P(gap$>$0)$=1.00$; P(gap$>$naive $0.03$)$=1.00$). The naive grader's $0.03$ had
hidden roughly seven-eighths of the true gap. All 200 human labels are released.

\paragraph{Takeaway.}
The measured register safety gap depends critically on how harm is judged, and the
danger runs opposite to intuition. A keyword refusal metric reports $0.27$--$0.54$; an
off-the-shelf LLM grader reports $0.03$, seeming to dismiss the gap as harmless
hedging; but human-validated grading reveals a real gap of $0.26$ ($95\%$ CI
$[0.19,0.32]$). The naive grader was not conservative---it was wrong, hiding most of
the harm. The refusal metric, though it counts hedging as unsafe, happened to track
the true harm magnitude closely here ($0.27$ vs $0.26$). The lesson is not that one
automatic metric is better, but that \textbf{an LLM safety grader must be validated
against human labels before its verdicts are believed}: used unchecked, it would have
reported no register safety problem where a substantial one exists.

\section{A Probe That Should Not Be Trusted}
\label{sec:probe}

To test whether models \emph{internally} represent the eval/deploy distinction, we
trained linear probes (logistic regression on last-token hidden states, 5-fold
cross-validated AUROC) to classify framing from activations at each layer. The
probe achieves AUROC $\approx 1.0$ at \emph{every} layer, including layer~1, for
every model (Figure~\ref{fig:probe}). This persists after (a) removing all
evaluation-marker vocabulary and (b) matching prompt length between framings.

We interpret this as a \emph{negative} methodological result rather than evidence of
internal evaluation-awareness. Perfect separation at the first layer---before
substantial contextual processing---is the signature of a surface shortcut: formal
and casual English differ in ordinary lexical distribution, and a linear probe
separates them trivially without reference to any latent ``being-tested'' concept.
Because a probe reading the prompt's own activations always has access to the
prompt's surface form, it cannot distinguish ``the model has inferred it is being
evaluated'' from ``these two texts are stylistically different.'' Our result
therefore corroborates \citet{devbunova2026} and cautions against treating
prompt-activation probe accuracy \emph{of this kind} as evidence of internal
evaluation awareness. We emphasize that this is a statement about last-token
prompt-activation probes as we implement them, not a claim that \emph{no} probe
could succeed. It motivates (i) behavioral measures such as RVG and (ii) probe
designs that read activations at a \emph{shared, identical} continuation appended to
both framings, so that any separability cannot come from surface form. We leave the
latter to future work.

\begin{figure}[t]
\centering
\includegraphics[width=0.8\linewidth]{figs/fig3_probe.png}
\caption{Layer-wise linear-probe AUROC for classifying eval vs.\ deploy framing
from hidden states. The probe saturates at AUROC $\approx 1.0$ from the first layer
for all models, even after removing evaluation keywords and matching prompt length.
We interpret this as evidence that the probe separates \emph{register} trivially,
not that the model internally represents being evaluated.}
\label{fig:probe}
\end{figure}

% ------------------------------------------------------------
\section{Discussion}

\paragraph{Implications for safety evaluation.}
Our results imply that a benchmark's \emph{register} is a hidden variable in its
safety measurements. Two benchmarks with identical harmful content but different
phrasing conventions can yield materially different refusal rates from the same
model. Because real deployment traffic is casual while most safety benchmarks are
formal, standard evaluations likely \emph{overestimate} deployed safety---by as much
as 54 percentage points of refusal rate in the case of Qwen-7B. Evaluators should
either match the register of their intended deployment or, at minimum, report
refusal rates under multiple registers.

\paragraph{Why cross-family evidence matters.}
Had we studied only Qwen, we would have reported a clean ``gap grows with scale''
result and a plausible emergent-capability narrative. Gemma refutes it. The
divergence is not noise but a consequence of different safety-tuning strategies:
Qwen preserves helpfulness under casual framing (so its gap grows as formal
refusals rise), while Gemma refuses aggressively in all registers (so its gap
shrinks toward a ceiling). This is a concrete example of why safety-evaluation
findings should be replicated across families before being treated as general.

\subsection{What ``register'' comprises: a threat to validity}
Our manipulation varies register as a \emph{bundle} of co-occurring stylistic
features: first-person framing, informal grammar and punctuation, politeness and
backstory, and discourse structure. We control for the two confounds most likely to
produce a trivial effect---explicit evaluation keywords and prompt
length---but we do not isolate which \emph{component} of register drives the
behavioral gap. It is therefore possible that the effect is attributable to a
specific sub-feature (e.g.\ first-person narration or the presence of a stated
motive) rather than to formality per se. We frame our claim accordingly: safety
behavior is sensitive to the stylistic difference between benchmark-like and
user-like phrasing, whichever component of that difference is responsible.
Disentangling the components---by ablating first-person framing, politeness, and
punctuation independently---is important future work and would sharpen the causal
interpretation. We flag this explicitly because it is, in our view, the most
important open question about the effect.

% ------------------------------------------------------------
\section{Limitations}
\label{sec:limitations}
\begin{itemize}
    \item \textbf{Limited intent diversity.} Although our dataset contains 300 paired
    prompts, these derive from only 25 underlying harmful intents. Our clustered
    bootstrap accounts for the resulting statistical dependence, but the limited
    diversity of underlying intents constrains how broadly the observed register effect
    can be generalized. Expanding the intent set while preserving the paired design is
    an important direction for future work.
    \item \textbf{Generation length cap.} Responses were capped at 200 new tokens.
    Actionable harmful content sometimes materializes only after a longer preamble, so
    a response truncated mid-buildup may be graded \textsc{hedge} when a longer
    generation would have been \textsc{harmful}. This deflates \textsc{harmful} counts
    and could do so asymmetrically if casual responses spend more tokens on backstory
    acknowledgment before content.
    \item \textbf{Self-grading.} Our harm grader is Qwen2.5-7B, which also generated
    part of the corpus. Self-preference is a known bias risk; a grader from a
    different family, or fuller human adjudication, would strengthen the harm-gap
    estimate.
    \item \textbf{Refusal $\neq$ harm.} Refusal-rate RVG measures whether a model
    refuses, not whether a non-refusal is harmful. We address this in
    Section~\ref{sec:harm} via harm grading; the residual limitation is that our
    corrected harm gap relies on audited grader-error rates from a modest validation
    sample (see below).
    \item \textbf{One casual style.} Our deploy framings share a single informal
    voice. If that voice has a stylistic fingerprint, we may be measuring sensitivity
    to \emph{that fingerprint} rather than to casualness in general. Multiple,
    independently written casual variants per intent (ideally crowdsourced) are needed
    to establish that the effect generalizes across casual styles; this is a more
    fundamental generalization question than the component-isolation issue discussed in our threats-to-validity subsection.
    \item \textbf{Quantization.} All models were run in 4-bit precision. Quantization
    can shift refusal behavior, so while the within-study comparison is internally
    consistent, absolute rates and possibly the gap may differ at full precision. A
    full-precision spot-check on the smallest models is a useful robustness step.
    \item \textbf{Single greedy generation.} We use one deterministic (greedy)
    generation per prompt, which captures no sampling variance. Refusal decisions can
    vary across seeds under temperature sampling, especially near a ceiling; our
    point estimates should therefore be read as greedy-decode measurements, and a
    sampled re-run of a subsample is a useful robustness check we leave to future
    work.
    \item \textbf{Refusal scoring.} We use a keyword/heuristic scorer, which can
    miss ``soft'' refusals that do not begin with canonical phrases. This most
    likely \emph{understates} eval-framing refusals, making our RVG a conservative
    estimate. Our audits (Section~\ref{sec:methods}) mitigate but do not eliminate
    this concern; a dedicated safety-judge classifier would sharpen the numbers.
    \item \textbf{Two families, small scales.} Six models across two families is
    cross-family evidence, not a claim about LLMs in general, and our largest model
    is 9B. The opposite scaling behaviors we observe are themselves a warning
    against over-generalizing from small samples.
    \item \textbf{Register is one axis.} Deployment framings are style-controlled
    rewrites, not verbatim user messages; real traffic differs in additional ways
    (topic distribution, multi-turn context) not captured here.
    \item \textbf{Probe scope.} Our negative probe result concerns last-token
    prompt-activation probes; it does not rule out that better-designed probes
    (e.g.\ shared-suffix or paired-decorrelated designs) could isolate a genuine
    representation.
\end{itemize}

% ------------------------------------------------------------
\section{Future Work}
Natural extensions include: replacing the keyword scorer with a safety-judge model
and reporting scorer agreement; extending to additional families and larger scales
(14B--70B) to map how far each family's trend continues; grounding deployment
framings in real user corpora rather than controlled rewrites; and designing
shared-suffix probes that read activations at an identical continuation to test for
genuine internal representations rather than surface separability.

% ------------------------------------------------------------
\section{Conclusion}
Instruction-tuned language models refuse casually-phrased harmful requests less often
than identical formally-phrased ones, and---contrary to what an off-the-shelf harm
grader would report---this corresponds to a real and substantial increase in
actionable harmful output. A keyword refusal metric reports a register gap of $0.27$
(up to $0.54$ for Qwen-7B); an unvalidated LLM grader reports only $0.03$; but
validating that grader against 200 human labels reveals it hides most casually-framed
harm as ``hedging,'' and the corrected register harm gap is $0.26$ ($95\%$ CI
$[0.19,0.32]$), with casual framing raising the actionable-harm rate from $8.5\%$ to
$34.3\%$. Our central conclusion is methodological and, we believe, broadly relevant
as LLM-graded evaluation proliferates: \textbf{an unvalidated LLM safety grader can
report the absence of a safety problem that in fact exists}, and must be checked
against human judgment before its verdicts are trusted. Within the two evaluated model
families, the register effect is consistently observed across all tested models; we
also find that prompt-activation probes
cannot validly establish internal evaluation-awareness for this distinction. Our study
is limited to open models up to 9B and a single grader; extending the human-validated
harm measurement to more models, graders, and scales is important future work.

\section*{Reproducibility Statement}
We release the paired-prompt dataset (all three iterations), the raw model
generations for all six models, the refusal labels, and the analysis code at
[repository link]. All experiments run on a single consumer GPU in 4-bit
precision, making the study reproducible without specialized hardware.

\section*{Ethics and Broader Impact}
Our prompts describe harmful intents in the abstract solely to measure refusal
behavior; we report results at the level of aggregate refusal rates and redact any
operationally harmful content from released generations. The register-sensitivity
gap we document is a safety-relevant weakness: it indicates that models may comply
with casually phrased harmful requests they would refuse under formal testing. We
disclose it to improve evaluation practice, not to provide a method for eliciting
harmful content; the effect is reported so that evaluators can correct for it.

% ------------------------------------------------------------
\begin{thebibliography}{99}
\small

\bibitem[Abdelnabi and Salem(2025)]{abdelnabi2025}
S.~Abdelnabi and A.~Salem.
\newblock The Hawthorne effect in reasoning models: Evaluating and steering test
awareness.
\newblock \emph{arXiv preprint arXiv:2505.14617}, 2025.

\bibitem[Arditi et al.(2024)]{arditi2024}
A.~Arditi, O.~Obeso, A.~Syed, D.~Paleka, N.~Panickssery, W.~Gurnee, and N.~Nanda.
\newblock Refusal in language models is mediated by a single direction.
\newblock \emph{arXiv preprint arXiv:2406.11717}, 2024.

\bibitem[Balepur et al.(2024)]{balepur2024}
N.~Balepur, A.~Ravichander, and R.~Rudinger.
\newblock Artifacts or abduction: How do LLMs answer multiple-choice questions
without the question?
\newblock In \emph{Proceedings of the 62nd Annual Meeting of the Association for
Computational Linguistics (Volume 1: Long Papers)}, pages 10308--10330, 2024.

\bibitem[Devbunova(2026)]{devbunova2026}
V.~Devbunova.
\newblock Is evaluation awareness just format sensitivity? Limitations of
probe-based evidence under controlled prompt structure.
\newblock \emph{arXiv preprint arXiv:2603.19426}, 2026.

\bibitem[Gemma Team(2024)]{gemma2}
Gemma Team.
\newblock Gemma 2: Improving open language models at a practical size.
\newblock Technical report, Google DeepMind. \emph{arXiv preprint
arXiv:2408.00118}, 2024.

\bibitem[Goldowsky-Dill et al.(2025)]{goldowsky2025}
N.~Goldowsky-Dill, B.~Chughtai, S.~Heimersheim, and M.~Hobbhahn.
\newblock Detecting strategic deception using linear probes.
\newblock \emph{arXiv preprint arXiv:2502.03407}, 2025.

\bibitem[Laine et al.(2024)]{laine2024}
R.~Laine et al.
\newblock The situational awareness dataset (SAD) for LLMs.
\newblock \emph{arXiv preprint arXiv:2407.04694}, 2024.

\bibitem[Mizrahi et al.(2024)]{mizrahi2024}
M.~Mizrahi, G.~Kaplan, D.~Malkin, R.~Dror, D.~Shahaf, and G.~Stanovsky.
\newblock State of what art? A call for multi-prompt LLM evaluation.
\newblock \emph{Transactions of the ACL}, 2024.

\bibitem[Needham et al.(2025)]{needham2025}
J.~Needham, G.~Edkins, G.~Pimpale, H.~Bartsch, and M.~Hobbhahn.
\newblock Large language models often know when they are being evaluated.
\newblock \emph{arXiv preprint arXiv:2505.23836}, 2025.

\bibitem[Nguyen et al.(2025)]{nguyen2025}
J.~Nguyen, K.~Hoang, C.~L.~Attubato, and F.~Hofst\"atter.
\newblock Probing evaluation awareness of language models.
\newblock \emph{arXiv preprint arXiv:2507.01786}; ICML Workshop on Technical AI
Governance, 2025.

\bibitem[Anonymous et al.(2026)]{notoneability2026}
Anonymous authors.
\newblock Evaluation awareness is not one capability: Evidence from open language
models.
\newblock \emph{arXiv preprint arXiv:2606.23583}, 2026.

\bibitem[Qwen Team(2025)]{qwen2025}
Qwen Team.
\newblock Qwen2.5 technical report.
\newblock Technical report, Alibaba Group. \emph{arXiv preprint arXiv:2412.15115},
2025.

\bibitem[Ren et al.(2024)]{ren2024safetywashing}
R.~Ren et al.
\newblock Safetywashing: Do AI safety benchmarks actually measure safety progress?
\newblock In \emph{Advances in Neural Information Processing Systems (NeurIPS),
Datasets and Benchmarks Track}. \emph{arXiv preprint arXiv:2407.21792}, 2024.

\bibitem[Xiong et al.(2025)]{xiong2025}
L.~Xiong, N.~Bhargava, J.~Hong, J.~Chang, H.~Liu, V.~Sharma, and K.~Zhu.
\newblock Probe-rewrite-evaluate: A workflow for reliable benchmarks and
quantifying evaluation awareness.
\newblock \emph{arXiv preprint arXiv:2509.00591}, 2025.

\bibitem[Zhou et al.(2023)]{zhou2023}
K.~Zhou et al.
\newblock Don't make your LLM an evaluation benchmark cheater.
\newblock \emph{arXiv preprint arXiv:2311.01964}, 2023.

\end{thebibliography}

\end{document}
